% ============================================================
%  EngineeringMath — English (LTR) preamble
%  Compile with XeLaTeX (or tectonic).  \EMroot must point to template/
% ============================================================
\input{\EMroot/em-packages}
\usepackage{fontspec}
\input{\EMroot/em-fonts-en}
\newfontfamily\EMfarsi[Path=\EMroot/fonts/, Extension=.ttf, Script=Arabic, AutoFakeBold=1.3]{Vazirmatn-Regular}
\newfontfamily\EMfaDisplay[Path=\EMroot/fonts/, Extension=.ttf, Script=Arabic]{Vazirmatn-Black}
\usepackage{bidi}
\input{\EMroot/em-math}

\linespread{1.08}
\newif\ifEMrtl \EMrtlfalse
% \rl{} for Persian snippets (Persian font + RTL)
\newcommand{\rl}[1]{\RL{\EMfarsi #1}}

% language strings
\def\EMlangCourse{Engineering Mathematics}
\def\EMlangCourseOther{\rl{ریاضی مهندسی}}
\def\EMlangUni{Ferdowsi University of Mashhad}
\def\EMlangUniOther{\rl{دانشگاه فردوسی مشهد}}
\def\EMlangChapter{Chapter}
\def\EMlangPart{Part}
\def\EMlangPage{Page}
\def\EMlangContents{Contents}
\def\EMlangChapterContents{In this chapter}
\def\EMlangFigure{Figure}
\def\EMlangDefinition{Definition}
\def\EMlangLemma{Lemma}
\def\EMlangProposition{Proposition}
\def\EMlangCorollary{Corollary}
\def\EMlangTheorem{Theorem}
\def\EMlangImportant{Key result}
\def\EMlangExample{Example}
\def\EMlangExercise{Exercise}
\def\EMlangRemark{Remark}
\def\EMlangProof{Proof}
\def\EMlangSolution{Solution}
\def\EMlangInstructorLabel{Instructor}
\def\EMlangInstructor{Abbas Ebrahimi Moghadam}
\def\EMlangCompilerLabel{Compiled and edited by}
\def\EMlangCompiler{Aidin Shekari}
\def\EMlangNotes{Lecture Notes}
\def\EMcompilerURL{https://shekari.me}
\def\EMlangColophonTitle{About these notes}


\def\EMlegendDefinition{A concept being defined}
\def\EMlegendDefinitionText{Definitions are set in indigo boxes and are numbered within the chapter.}
\def\EMlegendTheorem{A theorem or a formula to know}
\def\EMlegendTheoremText{Theorems and lemmas are set in rose boxes.}
\def\EMlegendImportant{A key result}
\def\EMlegendImportantText{The results worth finding again quickly are framed.}
\def\EMlegendExample{A worked example}
\def\EMlegendExampleText{Examples are green; the statement comes first.}
\def\EMlegendSolutionText{The solution follows inside the same box.}
\def\EMlegendExercise{An exercise}
\def\EMlegendExerciseText{Exercises are amber and are left to the reader.}
\def\EMlegendRemark{A remark}
\def\EMlegendRemarkText{Remarks and consequences have a thin grey rule and no number.}
\newcommand{\EMcolophontext}{%
  These are the typeset lecture notes of the course \emph{Engineering Mathematics},
  taught by \EMlangInstructor. They follow the instructor's slides and keep their order, notation
  and examples, in three parts: \textbf{Fourier analysis} (chapters 1--3), \textbf{partial differential
  equations} (chapters 4--7), and \textbf{complex analysis} (chapters 8--14).

  The notes were compiled and typeset by \EMlangCompiler. Every figure is drawn in TikZ or
  PGFPlots from the underlying mathematical object. Slips in the slides that were corrected
  while typesetting are listed in the file \texttt{REVIEW.md} of the repository.

  \medskip\noindent\textbf{How to read the notes.} Numbered boxes share one counter within each
  chapter, so "Example 5.12" is easy to find:}

\newcommand{\EMother}[1]{\rl{#1}}
\input{\EMroot/em-style}
